\documentclass{beamer}
\usetheme{Madrid}
\setbeamertemplate{navigation symbols}{}
\setbeamertemplate{footline}[frame number]
\usepackage{graphicx,amsmath,amssymb,booktabs,pgfpages,array}
\graphicspath{{../figures/}}
\newcommand{\FDMOrder}{1.9999}
\newcommand{\FEMOrder}{1.9990}
\expandafter\def\csname FD1\endcsname{393}
\expandafter\def\csname FD1Coarse\endcsname{44}
\expandafter\def\csname FE1\endcsname{441}
\expandafter\def\csname FE1Coarse\endcsname{53}
\expandafter\def\csname FD10\endcsname{943}
\expandafter\def\csname FD10Coarse\endcsname{90}
\expandafter\def\csname FE10\endcsname{945}
\expandafter\def\csname FE10Coarse\endcsname{94}
\expandafter\def\csname FD100\endcsname{2556}
\expandafter\def\csname FD100Coarse\endcsname{168}
\expandafter\def\csname FE100\endcsname{2568}
\expandafter\def\csname FE100Coarse\endcsname{189}
\newcommand{\ProfileN}{31}
\newcommand{\IterationN}{127}
\newcommand{\IterationCoarseN}{15}
\newcommand{\FDMSlice}{0.484}
\newcommand{\FEMSlice}{0.506}

\newcommand{\q}{\mathbf q}
\newcommand{\pic}[2][.95]{\begin{center}\includegraphics[width=#1\linewidth]{#2}\end{center}}
\title{Heat Flow Across a Material Interface}
\subtitle{Conservative FDM and piecewise-linear FEM}
\author{MTH321 Project 2}
\institute{XJTLU}
\date{November 2026}
\begin{document}
\begin{frame}
\titlepage
\centering What must a numerical method preserve when conductivity jumps?
\end{frame}

\begin{frame}{Two conductivity models}
\[-\nabla\cdot(a\nabla u)=f,\quad \q=-a\nabla u,\quad u=g\text{ on }\partial\Omega.\]
\pic[.85]{slides_conductivity.pdf}
\small On $\Omega=(0,1)^2$: smooth $a=1+x$ tests accuracy; layered
$a_L=1$, $a_R=10$ tests interface physics. Colors show conductivity;
the two color scales differ.
\end{frame}

\begin{frame}{Heat flux in the two models}
\pic{flux_vector_field.pdf}
\small Background: exact temperature. Arrows: exact $\q=-a\nabla u$.
Each panel has its own arrow key; compare magnitudes through those keys.
\medskip

The layered flux points left and is constant across the interface $x=1/2$.
\end{frame}

\begin{frame}{What the interface must preserve}
For perfect contact with no interface source,
\[[u]=0,\qquad [\q\cdot\mathbf n]=0.\]
The layered benchmark has $f=0$ and exact solution
\[u^*(x)=\begin{cases}20x/11,&x\leq1/2,\\(2x+9)/11,&x\geq1/2.\end{cases}\]
\[u_x^-=20/11,\quad u_x^+=2/11,\quad q_x^*=-20/11.\]
Temperature is continuous; its slope changes so that flux stays continuous.
\medskip

We impose $g=u^*$ on \textbf{all four boundary sides}.
\end{frame}

\begin{frame}{Two numerical representations}
\pic[.88]{methods_geometry.pdf}
\textbf{FDM}: balance fluxes shared by neighboring nodes.\\
\textbf{FEM}: approximate temperature by continuous affine functions on triangles.
\medskip

\small Left arrows indicate stencil directions. The black line on the right is the fitted material interface.
\end{frame}

\begin{frame}{FDM: coefficients and boundary assembly}
\[q_{i+1/2,j}=-a_{i+1/2,j}\frac{U_{i+1,j}-U_{i,j}}h.\]
Taking the discrete divergence gives
\[(A_hU)_P=\frac1{h^2}\sum_{s=E,W,N,S}a_s(U_P-U_s).\]
\begin{itemize}
\item Smooth coefficient: $a_{i+1/2,j}=(a_{i,j}+a_{i+1,j})/2$.
\item Share the same $a_s$ in the two adjacent rows.
\item For a boundary neighbor, $U_s=g_s$: add $a_sg_s/h^2$ to the right-hand side.
\end{itemize}
Positive shared weights and boundary elimination yield a symmetric positive-definite system.
\end{frame}

\begin{frame}{FEM: coefficients and boundary assembly}
For test functions $v_h=0$ on the boundary,
\[\int_\Omega a\nabla u_h\cdot\nabla v_h=\int_\Omega fv_h.\]
Write $u_h=\sum_j U_j\phi_j$ and assemble element contributions:
\[K_{ij}=\sum_T\int_T a\nabla\phi_j\cdot\nabla\phi_i,
\qquad F_i=\sum_T\int_T f\phi_i.\]
Separate interior and boundary degrees of freedom:
\[K_{II}U_I=F_I-K_{IB}g_B.\]
\begin{itemize}
\item Smooth $a$: quadrature integrates coefficient and load terms.
\item Layered $a$: fitted elements lie within one material each.
\end{itemize}
\small Positive weak energy gives SPD after Dirichlet elimination.
\end{frame}

\begin{frame}{How we measure error}
\textbf{Average error over the domain:}
\[E_{L^2}=\left(\int_\Omega|u^*-u_h|^2\,dx\right)^{1/2}.\]
\textbf{Largest error at the sampled nodes:}
\[E_{\infty,\mathrm{nodes}}=\max_i|u^*(x_i)-U_i|.\]
FDM: integrate a reconstructed function; FEM: integrate the P1 function.
The nodal maximum is a discrete $L^\infty$ measure, not the continuous supremum.
\medskip

For successive mesh sizes $h,\widetilde h$,
\[p=\frac{\log(E_h/E_{\widetilde h})}{\log(h/\widetilde h)}.\]
\end{frame}

\begin{frame}{Why smooth FDM is second order}
At a face midpoint $m$,
\[\frac{u(x+h)-u(x)}h=u'(m)+\frac{h^2}{24}u'''(m)+O(h^4).\]
The coefficient average also approximates $a(m)$ to order $h^2$.
Taking the flux difference gives the consistency defect
\[\tau_h=b_h-A_hR_hu^*=O(h^2).\]
The nodal error satisfies
\[A_h(U_h-R_hu^*)=\tau_h.\]
\textbf{Consistency + stability} gives second-order solution error for this smooth, uniformly elliptic problem.
\note{Smoothness is needed for the expansion and stability for the inference.
This argument must not be carried across a discontinuous coefficient.}
\end{frame}

\begin{frame}{Why P1 FEM is second order in $L^2$}
Let $B(w,v)=\int_\Omega a\nabla w\cdot\nabla v$ and $e=u-u_h$.
Galerkin orthogonality gives $B(e,v_h)=0$.
\[\|e\|_a\leq\|u-I_hu\|_a\leq Ch|u|_{H^2}.\]
For the dual problem $B(v,z)=(v,e)$,
\[\|e\|_{L^2}^2=B(e,z-I_hz)
 \leq Ch\|e\|_a\|z\|_{H^2}.\]
With dual regularity $\|z\|_{H^2}\leq C\|e\|_{L^2}$,
\[\boxed{\|e\|_{L^2}\leq Ch^2},\qquad \|e\|_a=O(h).\]
\small This does not by itself prove an unconditional $O(h^2)$ maximum-norm bound.
\end{frame}

\begin{frame}{Separating numerical error sources}
\small
\[\text{total error}=\text{discretization}+\text{quadrature}+\text{algebraic error}.\]
The smooth benchmark uses
\[u^*=\sin(\pi x)\sin(\pi y),\quad a=1+x,\quad g=0,\]
\[f=2\pi^2(1+x)\sin(\pi x)\sin(\pi y)-\pi\cos(\pi x)\sin(\pi y).\]
\begin{itemize}
\item Independently check $f$ by differentiating the analytic flux numerically.
\item Estimate algebraic error through a residual-correction solve.
\item Compare FEM quadrature orders 6 and 8.
\end{itemize}
Both extra effects are below $10^{-3}$ of the smooth measured $L^2$ error.
\note{Residual correction is a numerical estimate, not a certified roundoff bound.}
\end{frame}

\begin{frame}{Smooth benchmark: convergence in both measures}
\pic{smooth_convergence.pdf}
\small Four meshes: $N=16,32,64,128$, with $h=1/N$.
Left: continuous $L^2$; right: nodal maximum.
Dashed black lines indicate slope two; fitted orders use the finest three meshes.
\medskip

Both methods approach second order here; error constants differ.
\end{frame}

\begin{frame}{Why the interface needs a different argument}
At $x=1/2$, $a$ and $u_x$ jump. A Taylor expansion through that point is invalid.
\medskip

The physically useful invariant is instead
\[q_x=-a u_x=\text{constant}.\]
\begin{itemize}
\item \textbf{FDM}: what coefficient gives the correct resistance across an edge?
\item \textbf{FEM}: can the approximation space represent the slope change?
\end{itemize}
\medskip The two methods answer these questions differently.
\end{frame}

\begin{frame}{FDM: derive the interface resistance}
Integrate $u_x=-q_x/a$ along an edge crossing the interface:
\[u_R-u_L=-q_x\left(\frac{d_L}{a_L}+\frac{d_R}{a_R}\right).\]
Match this to $q_x=-a_{\rm edge}(u_R-u_L)/h$, where $h=d_L+d_R$:
\[\boxed{a_{\rm edge}=\frac{h}{d_L/a_L+d_R/a_R}}.\]
For a bisected edge,
\[a_H=\frac{2a_La_R}{a_L+a_R}=\frac{20}{11}.\]
The two material resistances add in series.
\end{frame}

\begin{frame}{Why arithmetic averaging fails here}
For the same edge,
\[a_A=\frac{1+10}{2}=5.5,\qquad a_H=\frac{20}{11}\approx1.818.\]
Inserting the \emph{exact nodal temperature drop} into each flux formula gives
\[\frac{q_A[u^*]}{q^*}=\frac{a_A}{a_H}=\frac{121}{40}.\]
\begin{itemize}
\item Arithmetic averaging assigns too little thermal resistance.
\item The computed temperature adjusts when the full system is solved.
\item $121/40$ is a consistency check, not the ratio of re-solved fluxes.
\end{itemize}
\small Positive coefficients can still give symmetry and boundedness while encoding the wrong interface physics.
\end{frame}

\begin{frame}{FEM: why fit the interface?}
On each triangle, P1 FEM has a constant gradient.
\begin{itemize}
\item A triangle crossing the interface cannot represent both slopes.
\item A fitted mesh places the kink on element edges.
\item Integrate $a$ within each material; impose exact boundary data.
\end{itemize}
For this benchmark, $u^*$ is continuous and affine within each material:
\[u^*\in V_h,\qquad B(u^*-u_h,v_h)=0.\]
Choosing $v_h=u^*-u_h$ gives
\[\|u^*-u_h\|_a^2=0.\]
Hence fitted FEM reproduces this solution in exact arithmetic.
\end{frame}

\begin{frame}{Three methods in the layered comparison}
\small
\begin{tabular}{>{\raggedright\arraybackslash}p{.26\linewidth}>{\raggedright\arraybackslash}p{.31\linewidth}>{\raggedright\arraybackslash}p{.31\linewidth}}
\toprule Method & Interface treatment & Plot identity\\\midrule
FDM arithmetic & Arithmetic face coefficient & Orange circles; solid line\\[1em]
FDM resistance & Subedge resistance; harmonic for equal lengths & Blue hollow squares; dashed line\\[1em]
FEM fitted & Interface on triangle edges; material-wise integration & Purple crosses; dotted line\\\bottomrule
\end{tabular}
\medskip

FDM uses odd $N=31,63,127,255$; FEM inserts the interface column.
Macro spacing is comparable, but unknown counts differ.
\end{frame}

\begin{frame}{Layered refinement: $L^2$ and nodal maximum error}
\pic{layered_convergence.pdf}
\small As for the smooth model: continuous $L^2$ error on the left,
nodal maximum error on the right, with $h=1/N$ on four meshes.
\medskip

Arithmetic errors decrease; the other two methods already have tiny errors
that do not follow the smooth model's convergence pattern. Why?
\end{frame}

\begin{frame}{Why the usual convergence test is inconclusive here}
This benchmark is \textbf{exactly representable}, rather than simply very smooth:
\begin{itemize}
\item $u^*$ is affine within each layer, with a kink at the interface.
\item Resistance FDM gives the exact nodal solution; fitted P1 FEM contains $u^*$.
\item Their discretization error is zero in exact arithmetic. Floating-point
      computation and the linear solve leave small measured errors.
\end{itemize}
We cannot infer a general convergence order from those tiny errors.
The benchmark instead tests whether the interface is represented correctly.
\note{The arithmetic method is not exact and its error decreases under refinement.
The FDM continuous L2 norm here uses resistance-coordinate reconstruction in crossing
cells, split at the interface; FEM uses the fitted P1 function. Ordinary bilinear
FDM reconstruction smears the kink and is recorded separately. No slope is fitted
to the numerical-error floor.}
\end{frame}

\begin{frame}{Where does the layered temperature error occur?}
\pic[1]{layered_error_maps.pdf}
\small Signed nodal error $u_h-u^*$ at $N=\ProfileN$, with one shared color scale.
Vertical dotted lines mark the interface; horizontal dashed lines mark the
slices used below.
\medskip

Arithmetic error concentrates around the interface and varies with $y$.
Resistance FDM and fitted FEM are essentially zero on this scale.
\end{frame}

\begin{frame}{Inspect the flux field before taking a slice}
\pic[1]{layered_numerical_flux_vectors.pdf}
\small Arrows show numerical $\q_h$, enlarged with the same scale for all methods;
lengths are not normalized. Dotted lines mark $x=1/2$.
\medskip

The exact field is $\q^*=(-20/11,0)$: horizontal and independent of $y$.
Arithmetic FDM introduces small changes in both magnitude and direction.
\note{The maps sample interpolated FDM face fluxes and FEM element gradients
away from the boundary and interface. Lengths are scaled uniformly by 1/25 in
domain coordinates. The true differences remain small; they are not amplified
relative to the base flow. The numerical field need not be one-dimensional.}
\end{frame}

\begin{frame}{Horizontal slices reveal the signed flux error}
\pic{slides_layered_result.pdf}
\small Along $y=y_0$, plot $u_h-u^*$ and the interface-normal component
$q_{x,h}(x,y_0)$. FDM: $y_0=\FDMSlice$; FEM: $y_0=\FEMSlice$; $N=\ProfileN$.
The gray reference is $q_x^*=-20/11$; the dotted vertical line is the interface.
\medskip

This slice checks normal flux; it does not replace the full 2D field.
Blue squares and purple crosses recover the reference. Arithmetic FDM does not.
\note{The exact solution is independent of y and has qy=0, which motivates the
horizontal slice. Arithmetic FDM can be y-dependent and have nonzero qy, so a
single qx curve cannot characterize all its errors. The preceding full-domain
error and vector maps supply that context. Odd FDM N has no y=0.5 row; actual
slice locations are labelled.}
\end{frame}

\begin{frame}{What validates the interface result?}
\begin{itemize}
\item Compare signed flux with the exact $-20/11$.
\item Reconstruct fluxes independently on the two material sides;
      sharing one face flux is not an independent continuity check.
\item All layered solutions remain between boundary extrema $0$ and $1$.
\item Positive diagonal and nonpositive off-diagonal entries support a
      discrete maximum principle on these meshes.
\end{itemize}
\medskip
Bounded temperature does \emph{not} prove accuracy: arithmetic FDM passes that check too.
\end{frame}

\begin{frame}{Refinement and contrast increase solver work}
\centering Unpreconditioned CG: refining $N=\IterationCoarseN\to\IterationN$
\medskip

\begin{tabular}{r rr}\toprule
Contrast & Resistance FDM & Fitted FEM\\\midrule
1 & $\csname FD1Coarse\endcsname\to\csname FD1\endcsname$ & $\csname FE1Coarse\endcsname\to\csname FE1\endcsname$\\
10 & $\csname FD10Coarse\endcsname\to\csname FD10\endcsname$ & $\csname FE10Coarse\endcsname\to\csname FE10\endcsname$\\
100 & $\csname FD100Coarse\endcsname\to\csname FD100\endcsname$ & $\csname FE100Coarse\endcsname\to\csname FE100\endcsname$\\\bottomrule
\end{tabular}
\medskip

\raggedright Entries are iteration counts, with zero initial guess and
$\|b-AU\|_2/\|b\|_2\leq10^{-10}$ in every case.
\medskip

Counts show increasing difficulty in the tested cases; they are not
condition-number estimates or wall-time rankings.
\end{frame}

\begin{frame}{Cost versus accuracy: what is measured?}
\[T=T_{\rm assembly}+T_{\rm direct\ solve}.\]
\begin{itemize}
\item One warm-up, then three timed repetitions per case; plot the median.
\item Exclude error diagnostics, plotting and file I/O from the timing.
\item Use verified errors and the same mesh cases as the refinement studies.
\item Compare these implementations; FEM and FDM assembly costs differ.
\end{itemize}
\medskip
\small In 2D, unknown count and sparse assembly/storage scale as $O(h^{-2})$.
Sparse direct factorization can scale as $O(h^{-3})$ with suitable ordering;
actual timings depend on implementation and hardware.
\end{frame}

\begin{frame}{Smooth model: cost versus accuracy}
\pic{smooth_cost_accuracy.pdf}
\small Horizontal axis: median assembly plus direct-solve time in seconds.
Vertical axes: continuous $L^2$ and nodal maximum errors.
Each marker is one mesh; connecting lines follow refinement.
\medskip

FDM attains lower errors at lower measured cost in these runs.
This describes these implementations and this benchmark.
\end{frame}

\begin{frame}{Layered model: cost versus accuracy}
\pic{layered_cost_accuracy.pdf}
\small Left: continuous $L^2$ error. Right: maximum relative signed-flux error
along the sampled slice, $\max|q_{x,h}-q_x^*|/|q_x^*|$.
Time includes assembly plus direct solve, using the same timing protocol.
\medskip

Refinement reduces the arithmetic error. Exact representations already reach
numerical precision; added work does not give a comparable reduction.
\end{frame}

\begin{frame}{Conclusions}
\begin{enumerate}
\item Both methods approach second-order smooth $L^2$ accuracy.
\medskip
\item Correct interface resistance or a fitted approximation space reproduces
      this layered benchmark; symmetry alone is insufficient.
\medskip
\item Flux diagnostics expose errors that temperature plots can hide.
\medskip
\item Refinement and contrast increase work; accuracy gains depend on what
      error remains to be reduced.
\end{enumerate}
\end{frame}

\begin{frame}{Coverage and remaining scope}
\small
\begin{tabular}{>{\raggedright\arraybackslash}p{.38\linewidth}>{\raggedright\arraybackslash}p{.55\linewidth}}\toprule
Problem-pack task & Evidence and presentation emphasis\\\midrule
Assembly, symmetry, SPD & Derivations shown; matrix/eigenvalue diagnostics saved\\[.5em]
Smooth errors and verification & Both norms shown; full residual tables and forcing checks summarized\\[.5em]
Interface and normal flux & Central comparison; side-flux mismatch values in saved tables\\[.5em]
Boundary extrema & Checked and explained through matrix signs\\[.5em]
Refinement and contrast & Iteration comparison shown; cost study added\\\bottomrule
\end{tabular}
\medskip

\textbf{Not studied:} unresolved FEM, curved interfaces, contact resistance,
AMG comparisons, or rigorous floating-point error bounds.
\note{These extensions are not all required tasks. FDM versus FEM is the depth
comparison. Source: Project 2 problem pack, Direction 3 and depth directions.
Code, data and figures are reproducible through run_all.py. AI assistance is
recorded in notes/audit_and_predictions.md; human review and individual contribution
records remain team responsibilities.}
\end{frame}
\end{document}
